\documentclass[11pt]{letter}
\usepackage[utf8]{inputenc}
\usepackage[margin=1in]{geometry}
\usepackage{hyperref}

\address{Dr. Sayuj Krishnan S, MBBS, DNB (Neurosurgery) \\ Consultant Neurosurgeon and Spine Surgeon \\ Yashoda Hospitals, Malakpet \\ Hyderabad, India \\ \texttt{dr.sayujkrishnan@gmail.com}}

\signature{Dr. Sayuj Krishnan S}

\begin{document}

\begin{letter}{The Editors \\ \textit{eLife}}

\opening{Dear Editors,}

I am writing to submit our manuscript, \textbf{``The Allometric Trap: Scaling Laws Predict When Growing Organisms Cannot Afford Their Own Shape,''} for consideration as a Research Article in \textit{eLife}.

This work identifies a previously unrecognized physical constraint on vertebrate morphogenesis: the \textbf{Allometric Trap}. We derive a dimensionless scaling parameter --- the Bio-Gravitational Number $\mathcal{B}_g = EI/(MgL^2)$ --- that predicts when an organism's passive structure becomes insufficient to resist gravity and active metabolic maintenance is required. Cross-species analysis reveals that humans uniquely occupy the deepest position in this trap ($\mathcal{B}_g \approx 0.01$), explaining why the human spine depends entirely on metabolic energy to maintain its shape.

\textbf{Key contributions:}
\begin{enumerate}
    \item \textbf{A universal scaling law.} The $L^4$ vs $L^2$ scaling mismatch between gravitational demand and metabolic supply creates a transient Energy Deficit Window during growth. This law is not species-specific; it is a physical consequence of life in gravitational fields, connecting to tree height limits, animal body-plan constraints, and microgravity physiology.

    \item \textbf{Molecular-resolution evidence.} AlphaFold v6 structural analysis of 23 key proteins reveals a 72\% demand--supply anisotropy gap ($p = 0.011$): mechanosensory proteins required for gravity sensing are structurally expensive, while their metabolic regulators are intrinsically disordered and fragile. This provides the first molecular-level evidence for why growth-spurt vulnerability exists.

    \item \textbf{A mechanistic explanation for scoliosis.} Our Information--Elasticity Coupling framework, implemented as a Cosserat rod model, reproduces clinical spinal angles from first principles and predicts scoliosis onset through a supercritical bifurcation at a critical spine length. This reclassifies Adolescent Idiopathic Scoliosis (affecting 3\% of children) from a genetic mystery to a predictable thermodynamic scaling failure.

    \item \textbf{Six falsifiable predictions} spanning HOX genetics (mouse micro-CT), microgravity physiology (astronaut MRI), clinical biomechanics (AIS progression cohorts), and developmental biology (zebrafish ciliary mutants), each with quantitative thresholds.
\end{enumerate}

This work bridges theoretical physics, structural biology, and clinical medicine. All simulation code is open-source and fully reproducible. I believe the interdisciplinary scope and the framework's ability to generate testable predictions across multiple biological scales make it well-suited for \textit{eLife}'s broad readership.

\closing{Sincerely,}

\end{letter}
\end{document}
